\section{安全、鲁棒性与上线策略}

课程主体不是安全专讲，但其“环境、奖励、评估”框架直接决定上线风险。Agent 的输出会触发真实动作，错误可跨步骤累积，还可能主动寻找奖励捷径；因此本节把课堂机制延伸到最小权限、动作确认、审计、回滚和人类接管，而不把安全缩成一句拒答提示。

\subsection{从训练安全到运行时安全}
安全控制必须沿着数据、策略、工具和运行环境分层布置，因为任何单点过滤都看不到完整动作链。训练阶段可以加入拒绝与恢复轨迹，推理阶段限制动作空间，运行时再以审计和人类确认兜底；三层共同决定系统是否能在未知输入下保持边界。
课程虽以训练为主，但其评估框架天然延伸到上线安全。Agent 的风险点包括越权工具调用、错误自动化放大、长链路偏航和不可解释决策。

\begin{importantbox}{安全策略分层}
\begin{enumerate}
    \item 训练层：加入反例轨迹和拒绝策略数据；
    \item 推理层：策略约束、动作白名单、敏感操作确认；
    \item 监控层：实时审计、异常回滚、人类接管阈值。
\end{enumerate}
\end{importantbox}

\subsection{鲁棒性评估应覆盖失败恢复}
多数 benchmark 聚焦最终结果，但生产环境更关心 ``失败后是否能自救''。课程关于 open-ended evaluation 的讨论可直接映射到该问题：需要在评估中显式纳入恢复路径质量。

\begin{knowledgebox}{鲁棒性观测指标}
\begin{itemize}
    \item 首次失败后的恢复成功率；
    \item 恢复所需额外步骤数；
    \item 恢复过程中新增风险动作比例；
    \item 恢复后结果是否满足业务阈值。
\end{itemize}
\end{knowledgebox}

\subsection{分阶段上线}
对 Agent 产品，课程思想对应的上线方法是 ``先可控，再扩域''：先在闭集任务上线高置信动作，再逐步引入开集任务，并持续校准评估器。

\begin{warningbox}{上线反模式}
一次性放开所有工具权限、把开集任务直接按闭集规则打分、缺少失败复盘机制，这三种做法在 Agent 场景会显著放大风险与成本。
\end{warningbox}

\subsection{本章小结}
Agent 安全不是额外模块，而是训练、评估、推理、运维的联合产物。鲁棒性要通过失败恢复能力来度量，而非只看一次成功率。


